\documentclass[10pt]{article}
\usepackage[margin=1in]{geometry}
\usepackage[T1]{fontenc}
\usepackage{lmodern,microtype,amsmath,amssymb,amsthm}
\theoremstyle{definition}
\newtheorem{assumption}{Assumption}
\begin{document}
\section*{Restarted Schedule Gradient Descent}
We consider finding a stationary point of a possibly nonconvex function $f:E\to\mathbb{R}$ on a finite-dimensional real inner product space $E$. The function is twice continuously differentiable, with $L$-Lipschitz gradient and $\rho$-Lipschitz Hessian, where $L,\rho>0$. A given starting point $x^0\in E$ satisfies $f(x^0)-\inf_{x\in E}f(x)\leq\Delta$ for $\Delta\geq0$. The oracle returns $\nabla f(x)$ at $x$. For $\varepsilon>0$, restarted schedule gradient descent must return $\bar{x}\in E$ with $\|\nabla f(\bar{x})\|\leq\varepsilon$ for every realization.

The underlying optimization problem is
\[
\inf_{x\in E} f(x).
\]
The norm $\|\cdot\|$ is induced by the inner product on $E$, and $\|\cdot\|_{\mathrm{op}}$ denotes the induced operator norm. The algorithm receives $E,x^0,L,\rho,\Delta$, and $\varepsilon$.

Our assumptions:
\begin{assumption}[Gradient and Hessian regularity]
The function $f$ belongs to $C^2(E;\mathbb{R})$. Its Hessian $\nabla^2 f(x)$ is a linear operator from $E$ to $E$, and for all $x,y\in E$,
\[
\begin{aligned}
\|\nabla f(x)-\nabla f(y)\|&\leq L\|x-y\|,\\
\|\nabla^2 f(x)-\nabla^2 f(y)\|_{\mathrm{op}}&\leq\rho\|x-y\|.
\end{aligned}
\]
\end{assumption}

\begin{assumption}[Initial objective gap]
The given starting point satisfies
\[
f(x^0)-f(x)\leq\Delta,\qquad x\in E.
\]
\end{assumption}

Our result is algorithm-specific: we analyze restarted schedule gradient descent. The method chooses thresholds $G_1,R\in\mathbb{R}$ and a finite schedule $(h_0,\ldots,h_{m-1})$, with $m\geq0$ and $h_j>0$ for $0\leq j<m$, using the available inputs. Neither threshold has a sign restriction, and the schedule may be empty. Define $\widetilde h_j=h_j$ for $j<m$ and $\widetilde h_j=0$ for $j\geq m$.

The method maintains a point $x^k$, a restart centre $c^k$, and an index $j_k$, initialized by
\[
(x^0,c^0,j_0)=(x^0,x^0,0).
\]
At call $k$, it queries $g^k=\nabla f(x^k)$. If $\|g^k\|\leq\varepsilon$, it freezes the point, centre, and index for all subsequent calls. Otherwise, while the method has not frozen, it applies the following update. If $j_k=0$ and $\|g^k\|>G_1$, then
\[
\begin{aligned}
x^{k+1}&=x^k-L^{-1}g^k,\\
c^{k+1}&=c^k,\qquad j_{k+1}=0.
\end{aligned}
\]
In every other case, it sets
\[
\begin{aligned}
y^k&=\begin{cases}
x^k,&j_k=0,\\
c^k,&j_k>0,
\end{cases}\\
x^{k+1}&=x^k-\widetilde h_{j_k}g^k,\\
c^{k+1}&=y^k,\\
j_{k+1}&=\begin{cases}
j_k+1,&j_k+1<m\ \text{and }\|x^{k+1}-y^k\|\leq R,\\
0,&\text{otherwise}.
\end{cases}
\end{aligned}
\]
Queries lie in $E$. After an allotted number $N\geq0$ of oracle calls, the output is the current point $\bar{x}=x^N$. The required accuracy is
\[
\|\nabla f(\bar{x})\|\leq\varepsilon
\qquad\text{for every realization}.
\]
\end{document}
